\section{연습문제}
\label{sec:ring-exercises}

문제는 세 단계로 나뉜다. A는 정의와 계산, B는 핵심 증명, C는 구조를 연결하는 심화 문제이다. 별표가 붙은 문제는 뒤 절에 상세 풀이가 있다.

\subsection*{A. 정의와 계산}

\begin{exercise}
환 $R$에서 다음을 직접 증명하라.
\[
  a(-b)=-(ab),\qquad (-a)(-b)=ab,
  \qquad a(b-c)=ab-ac.
\]
\end{exercise}

\begin{exercise}
다음 부분집합이 주어진 환의 부분환인지, 아이디얼인지 각각 판단하라.
\begin{enumerate}[label=(\alph*)]
  \item $2\Z\subseteq\Z$.
  \item $\Z\subseteq\Q$.
  \item $\R\times\{0\}\subseteq\R\times\R$.
  \item $\{f\in\R[X]:f(0)=0\}\subseteq\R[X]$.
\end{enumerate}
\end{exercise}

\begin{exercise}[*]
$\Z/20\Z$의 모든 단위원과 모든 영이 아닌 영인수를 구하라. 각 단위원의 역원도 구하라.
\end{exercise}

\begin{exercise}
$\Z/24\Z$의 멱영원소를 모두 구하라.
\end{exercise}

\begin{exercise}
$R,S$를 영환이 아닌 가환환이라 하자. 환 $R\times S$의 단위원과 멱영원소를 각 성분의 자료로 기술하라. 또한 영이 아닌 $(r,s)$가 영인수일 필요충분조건을 $r,s$의 소멸원소를 이용하여 설명하라.
\end{exercise}

\begin{exercise}[*]
다음 아이디얼의 생성원을 가능한 한 단순하게 구하라.
\[
 (18,30)\subseteq\Z,\qquad
 (12)\cap(18)\subseteq\Z,\qquad
 (X^2,X^5)\subseteq K[X].
\]
\end{exercise}

\begin{exercise}
$\Z[X]$에서 다음 아이디얼을 원소의 계수 조건으로 설명하라.
\[
  (X),\qquad (3),\qquad (3,X),\qquad (X^2).
\]
\end{exercise}

\begin{exercise}[*]
평가 준동형
\[
  \operatorname{ev}_2:\Z[X]\to\Z,\qquad f\mapsto f(2)
\]
의 핵과 상을 구하고, 제1동형정리로 몫환을 기술하라.
\end{exercise}

\begin{exercise}
다음 몫환의 원소를 표준형으로 쓰고, 정역 또는 체인지 판단하라.
\begin{enumerate}[label=(\alph*)]
  \item $\Q[X]/(X^2-3)$.
  \item $\R[X]/(X^2-4)$.
  \item $\F_2[X]/(X^2+X+1)$.
  \item $\F_2[X]/(X^2+1)$.
\end{enumerate}
\end{exercise}

\begin{exercise}[*]
$\F_3[X]/(X^2+1)$의 모든 원소를 쓰고, 각 영이 아닌 원소의 역원을 구하라.
\end{exercise}

\subsection*{B. 핵심 증명}

\begin{exercise}
환의 아이디얼들의 임의의 교집합은 아이디얼임을 증명하라. 두 아이디얼의 합집합은 언제 아이디얼이 되는가?
\end{exercise}

\begin{exercise}[*]
환 $R$의 멱영원소 전체
\[
  \sqrt{(0)}=\{a\in R:a^n=0\text{인 }n\ge1\text{이 존재}\}
\]
가 아이디얼임을 증명하라.
\end{exercise}

\begin{exercise}
가환환 $R$의 진아이디얼 $I$에 대해 다음이 동치임을 보이라.
\begin{enumerate}[label=(\roman*)]
  \item $I$는 극대아이디얼이다.
  \item 모든 $a\notin I$에 대해 어떤 $r\in R$, $i\in I$가 존재하여 $ra+i=1$이다.
\end{enumerate}
\end{exercise}

\begin{exercise}[*]
환 준동형 $\varphi:R\to S$에 대해 다음을 증명하라.
\begin{enumerate}[label=(\alph*)]
  \item $J\trianglelefteq S$이면 $\varphi^{-1}(J)\trianglelefteq R$이다.
  \item $J$가 소아이디얼이면 $\varphi^{-1}(J)$도 소아이디얼이다.
  \item $\varphi$가 전사이고 $J$가 극대아이디얼이면 $\varphi^{-1}(J)$도 극대아이디얼이다.
\end{enumerate}
\end{exercise}

\begin{exercise}
환 준동형 $\varphi:R\to S$가 전사라고 하자. $I\trianglelefteq R$이면 $\varphi(I)$가 $S$의 아이디얼임을 증명하라. 전사 조건을 제거하면 거짓임을 예로 보여라.
\end{exercise}

\begin{exercise}[*]
$S\subseteq R$를 $R$과 같은 항등원을 갖는 부분환이라 하고 $J\trianglelefteq R$라 하자. $S+J$가 $R$의 부분환이고 $S\cap J$가 $S$의 아이디얼임을 확인한 뒤,
\[
  S/(S\cap J)\cong(S+J)/J
\]
를 환의 제1동형정리로 증명하라.
\end{exercise}

\begin{exercise}
$I,J\trianglelefteq R$이고 $I\subseteq J$라 하자. 동형
\[
  (R/I)/(J/I)\cong R/J
\]
를 대표원을 사용하여 직접 구성하고, 잘 정의됨을 확인하라.
\end{exercise}

\begin{exercise}[*]
$I+J=R$이면 $I^m+J^n=R$임을 모든 양의 정수 $m,n$에 대해 증명하라.
\end{exercise}

\begin{exercise}
유한 가환환 $R$의 모든 원소가 영인수 또는 단위원임을 증명하라. 이 문제에서만 $0$도 영인수로 허용한다고 하자.
\end{exercise}

\begin{exercise}[*]
$R\ne0$인 가환환에 대해 다음이 동치임을 증명하라.
\begin{enumerate}[label=(\roman*)]
  \item $R$은 체이다.
  \item $S\ne0$인 모든 환 $S$와 모든 환 준동형 $R\to S$에 대해 그 준동형은 단사이다.
  \item $R$의 아이디얼은 서로 다른 두 아이디얼 $(0)$과 $R$뿐이다.
\end{enumerate}
\end{exercise}

\subsection*{C. 구조와 응용}

\begin{exercise}[*]
중국인의 나머지 정리를 이용하여
\[
  x\equiv4\pmod9,\qquad
  x\equiv3\pmod{10},\qquad
  x\equiv5\pmod7
\]
을 풀라.
\end{exercise}

\begin{exercise}
양의 정수 $n=p_1^{e_1}\cdots p_r^{e_r}$가 소인수분해라고 하자. 환 동형
\[
  \Z/n\Z\cong
  \prod_{i=1}^r\Z/p_i^{e_i}\Z
\]
를 증명하라. 이로부터 단위원군의 원소 개수에 대한 오일러 곱공식을 유도하라.
\end{exercise}

\begin{exercise}[*]
$K$를 표수가 $2$가 아닌 체라 하자. 다음 동형을 구체적으로 구성하라.
\[
  K[X]/(X^2-1)\cong K\times K.
\]
$X+(X^2-1)$은 오른쪽에서 어떤 원소에 대응하는가?
\end{exercise}

\begin{exercise}
$K[X]/(X^3-X)$를 $K$의 표수에 따라 중국인의 나머지 정리로 가능한 한 많이 분해하라.
\end{exercise}

\begin{exercise}[*]
$A=\Q[X]/(X^3-2)$이고 $\alpha=X+(X^3-2)$라 하자. 다음을 계산하라.
\begin{enumerate}[label=(\alph*)]
  \item $(1+\alpha+\alpha^2)(1-\alpha)$.
  \item $(\alpha^2+\alpha+1)^{-1}$.
  \item $A$가 체인 이유.
\end{enumerate}
\end{exercise}

\begin{exercise}
$R=K[X,Y]/(XY)$라 하고 $x,y$를 $X,Y$의 잉여류라 하자.
\begin{enumerate}[label=(\alph*)]
  \item $x,y$가 영이 아님을 보이라.
  \item $xy=0$을 확인하라.
  \item $(x)$와 $(y)$가 소아이디얼임을 보이라.
  \item $(x,y)$가 극대아이디얼임을 보이라.
\end{enumerate}
\end{exercise}

\begin{exercise}[*]
$R$이 PID이고 $a,b\in R$라 하자. $(a,b)=(d)$이면 $d$가 $a,b$의 최대공약수임을 증명하고, 역으로 최대공약수 $d$가 항상 $(a,b)$를 생성하는지는 일반 UFD에서도 성립하는지 조사하라.
\end{exercise}

\begin{exercise}[*]
가우스 정수환 $\Z[i]$에서
\[
  5=(2+i)(2-i)
\]
임을 확인하고, 노름을 이용하여 $2+i$가 단위원이 아님을 보이라. 이어서 환 동형
\[
  \Z[i]/(2+i)\cong\F_5
\]
를 구체적으로 구성하고, 이 동형 아래에서 $i$가 $\F_5$의 어떤 원소에 대응하는지 구하라. 그 원소가 $-1$의 제곱근임도 확인하라.
\end{exercise}

\begin{exercise}[*]
$R$을 정역이라 하자.
\begin{enumerate}[label=(\alph*)]
  \item $R$의 모든 소아이디얼이 극대아이디얼이면 $R$이 체임을 증명하라.
  \item $R$이 PID이면 모든 영이 아닌 소아이디얼이 극대아이디얼임을 증명하라.
  \item $\Z$가 (b)는 만족하지만 (a)의 가정은 만족하지 않는 이유를 설명하라.
\end{enumerate}
\end{exercise}

\begin{exercise}
체 $K$와 기약다항식 $f\in K[X]$에 대해 $L=K[X]/(f)$라 하자. $\deg f=n$이면 $L$이 $K$-벡터공간으로 차원 $n$임을 보이라. $1,\alpha,\dots,\alpha^{n-1}$이 기저임을 직접 증명하라.
\end{exercise}
